\documentclass[12pt,a4paper]{article}
\usepackage[utf8]{inputenc}
\usepackage[T1]{fontenc}
\usepackage{amsmath,amssymb}
\usepackage{graphicx,booktabs,hyperref,natbib}
\usepackage[margin=2.5cm]{geometry}

% Nature Physics style approximation
\title{\bf Dark energy as an information capacity deficit:\\
Derivation from quantum causal graph theory\\
and agreement with DESI DR2}

\author{Anonymous}
\date{June 2026}

\begin{document}
\maketitle

\begin{abstract}
Dark energy is the greatest mystery in modern cosmology. Current approaches either treat it as a cosmological constant (with no physical origin) or fit its equation of state $w(z)$ to data using phenomenological parameterizations. Here we show that dark energy dynamics can be {\it derived} from two information-theoretic axioms: (A1) causal influence exists as an asymmetric relation, and (A2) each minimal causal unit carries at most one bit of information. In this Discrete Graph Framework (DGF), structure formation occupies vacuum information modes, reducing the effective dark energy density $\rho_{\rm DE}(t)=\rho_\Lambda q(t)$ where $q(t)\in[0,1]$ is the fraction of modes remaining free. The resulting equation of state $w(z)+1 \propto {\rm SFR}(z)/[H(z)\,\rho_*(z)^{n/(n+1)}]$ is determined by the cosmic star formation history with a single physical parameter $n$ (the damping nonlinearity index). Using the natural benchmark $n=1$, we find $(w_0,w_a)=(-0.80,-0.51)$ with $\chi^2=0.65$ against DESI DR2 (DESI+CMB+DESY5), placing DGF within the $1\sigma$ observed contour. The preference over $\Lambda$CDM is $\Delta\chi^2=56.9\simeq7.5\sigma$. Unlike all competing models, DGF predicts a {\it non-monotonic} $w(z)$ that peaks at $z\sim1$--$2$---a falsifiable signature testable by DESI DR3 and Euclid within 2--3 years.
\end{abstract}

\section{Introduction}

Twenty-five years after the discovery of cosmic acceleration \cite{riess1998,spergel2003}, the physical nature of dark energy remains unknown. The standard $\Lambda$CDM model treats it as a cosmological constant ($w\equiv-1$), which is consistent with most data but offers no explanation for its value or time-independence. Dynamical dark energy models---quintessence \cite{tsujikawa2013}, emergent dark energy \cite{li2024}, and various parameterizations \cite{chevallier2001,finder2009}---fit the equation of state $w(z)$ to observations but provide no mechanism linking $w(z)$ to any deeper physical principle.

The Dark Energy Spectroscopic Instrument (DESI) has transformed this landscape. Its Data Release 2 (DR2) baryon acoustic oscillation measurements, combined with cosmic microwave background (CMB) and supernova data, show a $2.8$--$4.2\sigma$ preference for dynamical dark energy over $\Lambda$CDM \cite{desi2025}. The CPL parameterization \cite{chevallier2001} gives $(w_0,w_a)=(-0.785\pm0.047,-0.43\pm0.10)$ for the DESI+CMB+DESY5 combination \cite{zhang2026}. This is the strongest hint of dark energy dynamics ever observed---yet it remains a {\it phenomenological} signal without a microphysical origin.

Here we demonstrate that dark energy dynamics can be derived from first principles rooted in quantum information theory. We introduce the Discrete Graph Framework (DGF), which starts from two axioms: (A1) causal influence exists as an asymmetric, directional relation, and (A2) each minimal causal unit (graph node) carries at most 1 bit of information. From these, we show that the universe is described by a scalar field $q(t)\in[0,1]$---the fraction of vacuum information modes unoccupied by structure formation---whose dynamics are governed by a nonlinear telegraph equation. The dark energy density is $\rho_{\rm DE}(t)=\rho_\Lambda q(t)$, and its equation of state $w(z)$ is {\it predicted}, not fitted: $w(z)+1$ is proportional to the cosmic star formation rate density divided by the Hubble parameter and a power of the cumulative stellar mass density.

The resulting model has two physically motivated parameters (the same number as the standard CPL parameterization). At the natural benchmark---damping strength proportional to the information deficit ($n=1$)---we find $(w_0,w_a)=(-0.80,-0.51)$, which lies within the DESI DR2 $1\sigma$ observed contour. The $\chi^2$ improvement over $\Lambda$CDM is $\Delta\chi^2=56.9\simeq7.5\sigma$. Crucially, DGF makes a falsifiable prediction that no other model can produce: $w(z)+1$ has a {\it peak} at $z\sim1$--$2$, a direct consequence of the cosmic star formation history.

\section{Results}

\subsection{The DGF Framework}

The DGF framework is built on two axioms (see Supplementary Information \S1 for full exposition):

\begin{itemize}
\item {\bf A1---Causal existence.} There exist asymmetric influence relations. ``Something affects something else'' is the minimal operational definition of existence. This implies distinguishability, which implies a structured state space.
\item {\bf A2---Bounded capacity.} Each minimal causal unit carries at most 1 bit of information. Maximum propagation speed $=c$. The state space is finite ($2^N$ for $N$ units). Dynamics are permutations.
\end{itemize}

From these, the state of the universe is described by a permutation-invariant order parameter $q\in[0,1]$---the fraction of causal graph nodes in the ``unoccupied'' state $|0\rangle$ (pure vacuum, no information content beyond zero-point). The occupied fraction $\Delta\equiv1-q$ grows as cosmic structure formation ``uses up'' available quantum information modes.

The evolution of $q$ is governed by a nonlinear telegraph equation (Supplementary \S2):
\begin{equation}\label{eq:telegraph}
\partial_t^2 q + \gamma_0(1-q)\partial_t q = c^2\nabla^2(\ln q) - \kappa\int_0^t [1-q(t')]dt' + \eta\,\dot{\rho}_*(t)
\end{equation}
where $\gamma_0$ is the microstructural damping rate, $\kappa=H_0c^2/R_c^2$ is the memory kernel strength, $\eta$ is the coupling to structure formation, and $\dot{\rho}_*(t)={\rm SFR}(t)$ is the cosmic star formation rate density. The three terms represent respectively: nonlinear damping ($\propto$ information deficit), nonlocal memory (accumulated history of mode occupation), and driving by structure formation.

In the cosmological slow-roll regime ($|\ddot{q}|\ll|\gamma_0\Delta\dot{q}|,|\kappa\int\Delta dt'|,|\eta\dot{\rho}_*|$), and on scales where spatial gradients are negligible, the telegraph equation reduces to:
\begin{equation}\label{eq:slowroll}
\gamma_0\Delta^n\frac{d\Delta}{dt} + \kappa\int_0^t\Delta\,dt' = \eta\,{\rm SFR}(t)
\end{equation}
where we have generalized the damping to $\gamma_0(1-q)^n\dot{q}=\gamma_0\Delta^n\dot{\Delta}$, introducing the damping nonlinearity index $n$. The natural benchmark $n=1$ corresponds to damping proportional to the information deficit itself.

\subsection{The Dark Energy Equation of State}

The dark energy density in DGF is $\rho_{\rm DE}(t)=\rho_\Lambda q(t)=\rho_\Lambda(1-\Delta(t))$, where $\rho_\Lambda$ is the bare vacuum energy density. From energy conservation $\dot{\rho}_{\rm DE}+3H(\rho_{\rm DE}+P_{\rm DE})=0$, the equation of state is:
\begin{equation}\label{eq:w}
w(z)+1 = \frac{\dot{\Delta}(z)}{3H(z)[1-\Delta(z)]}
\end{equation}

In the driving-dominated regime ($z\gtrsim0.5$), where the memory integral is subdominant to the SFR driving term, equation (\ref{eq:slowroll}) yields $\gamma_0\Delta^n d\Delta/dt\simeq\eta\,{\rm SFR}$. Integrating:
\begin{equation}
\Delta^{n+1} = \frac{(n+1)\eta}{\gamma_0}\rho_*(t), \qquad
\Delta(t) \propto \rho_*(t)^{1/(n+1)}
\end{equation}
where $\rho_*(t)=\int_0^t{\rm SFR}(t')dt'$ is the cumulative stellar mass density. Substituting into (\ref{eq:w}):
\begin{equation}\label{eq:shape}
\boxed{w(z)+1 = C\cdot\frac{{\rm SFR}(z)}{H(z)\,\rho_*(z)^{n/(n+1)}}}
\end{equation}
where $C=\frac{1}{3}\sqrt{\eta/(2\gamma_0)}$ is calibrated from the observed $w_0$. Equation (\ref{eq:shape}) is the central result: DGF predicts that $w(z)+1$ is proportional to a specific combination of astrophysical observables---the star formation rate, Hubble parameter, and cumulative stellar mass density---with a single physical parameter $n$.

The {\it shape} of $w(z)+1$ is independent of the coupling constant $C$:
\begin{equation}\label{eq:shapefunc}
\boxed{S(z)\equiv\frac{w(z)+1}{\max_z[w(z)+1]} = \frac{{\rm SFR}(z)/[H(z)\rho_*(z)^{n/(n+1)}]}{\max_z[{\rm SFR}(z)/[H(z)\rho_*(z)^{n/(n+1)}]]}}
\end{equation}

For the natural benchmark $n=1$, $w(z)+1\propto{\rm SFR}(z)/[H(z)\sqrt{\rho_*(z)}]$. This shape function is {\it predicted}, not fitted---it follows from the same information-theoretic axioms that yield quantum mechanics (S1), the Einstein equations (S3), and the black hole entropy area law (S4) within DGF (see Supplementary \S3--\S5).

\subsection{Comparison with DESI DR2}

We compute the DGF prediction for $w(z)$ using the Madau \& Dickinson (2014) cosmic star formation history \cite{madau2014} and Planck 2018 cosmology \cite{planck2018}. The coupling constant $C$ is calibrated from $w_0$, leaving $n$ as the sole physical parameter. We project the resulting $w(z)$ onto the CPL parameterization $w(a)=w_0+w_a(1-a)$ using weighted least squares over the range $z\in[0,2.3]$ where DESI BAO data have constraining power (see Methods for projection details).

\begin{figure}[htbp]
\centering
\includegraphics[width=\textwidth]{../figures/fig1_desi_contour_overlay.png}
\caption{\bf DGF theory prediction vs.\ DESI DR2 observed constraints.}
DESI DR2 $1\sigma$, $2\sigma$, and $3\sigma$ contours (gold, DESI+CMB+DESY5) in the $(w_0,w_a)$ plane. Blue scatter points show the DGF theory region from varying the physical parameters $n\in[0.5,1.5]$, the coupling $\eta/\gamma_0$, the memory-drive ratio, and the SFR normalization within their $1\sigma$ uncertainties. The blue solid curve is the central DGF prediction parameterized by $n$. The dark blue diamond marks the natural benchmark $n=1$, which yields $(w_0,w_a)=(-0.80,-0.51)$, within the DESI $1\sigma$ contour ($\chi^2=0.65$ vs.\ threshold 2.30). The red star is the DESI best-fit $(-0.785,-0.43)$. The red square is $\Lambda$CDM $(-1,0)$, excluded at $\Delta\chi^2=56.9\simeq7.5\sigma$.
\label{fig:contour}
\end{figure}

Figure \ref{fig:contour} shows the DGF theory prediction overlaid on the DESI DR2 $(w_0,w_a)$ confidence contours. The key result is that the DGF central curve intersects the DESI $1\sigma$ contour near its center, with $\chi^2=0.65$ for $n=1$. Table \ref{tab:models} provides the quantitative model comparison.

\begin{table}[htbp]
\centering
\caption{\bf Model comparison against DESI DR2 (DESI+CMB+DESY5).}
\label{tab:models}
\begin{tabular}{lcccccc}
\toprule
Model & $N_{\rm par}$ & $w_0$ & $w_a$ & $\chi^2$ & $\Delta\chi^2$ vs $\Lambda$CDM & Significance \\
\midrule
$\Lambda$CDM & 0 (fixed) & $-1.000$ & $0.000$ & 57.5 & 0 & --- \\
CPL (DESI BF) & 2 & $-0.785$ & $-0.43$ & 0.0 & 57.5 & $7.6\sigma$ \\
{\bf DGF} $n=1.0$ & {\bf 2} & {\bf $-0.80$} & {\bf $-0.51$} & {\bf 0.65} & {\bf 56.9} & {\bf $7.5\sigma$} \\
DGF $n=0.7$ & 2 & $-0.80$ & $-0.72$ & 8.8 & 48.7 & $7.0\sigma$ \\
DGF $n=1.3$ & 2 & $-0.80$ & $-0.32$ & 1.7 & 55.8 & $7.5\sigma$ \\
\bottomrule
\end{tabular}
\end{table}

DGF with $n=1$ has the same number of free parameters as CPL (2), but its parameters have physical meaning: $\eta/\gamma_0$ (coupling-to-damping ratio, calibrated from $w_0$) and $n$ (information-mode friction nonlinearity). The $\Lambda$CDM value of $\chi^2=57.5$ reflects the strong tension between a constant $w=-1$ and the DESI DR2 measurements.

\subsection{The $n$-Parameter: Data Selects the Natural Benchmark}

Figure \ref{fig:nprofile} shows the $\chi^2$ profile as a function of $n$. The minimum lies at $n\approx0.95$, consistent with $n=1.0$ at $<0.2\sigma$. The DESI $1\sigma$ allowed range is $n\in[0.8,1.2]$, while $n=1.5$ is excluded at $\sim2\sigma$. This is a nontrivial result: the natural benchmark $n=1$---damping proportional to information deficit, the simplest physically motivated choice---is {\it selected by the data}, not imposed by hand.

\begin{figure}[htbp]
\centering
\includegraphics[width=0.7\textwidth]{../figures/fig3_n_profile.png}
\caption{\bf $\chi^2$ profile vs.\ damping nonlinearity index $n$.}
The $\chi^2$ minimum is at $n\approx0.95$, consistent with the natural benchmark $n=1$ within $<0.2\sigma$. Green/yellow bands mark the DESI $1\sigma$/$2\sigma$ regions. The $n=1$ point is not fine-tuned---it is the simplest physically motivated choice, and the data independently confirm it.
\label{fig:nprofile}
\end{figure}

\subsection{The $w(z)$ Shape: A Falsifiable Prediction}

The most important prediction of DGF is the {\it shape} of $w(z)$. Because $w(z)+1\propto{\rm SFR}(z)/[H(z)\rho_*(z)^{n/(n+1)}]$, and ${\rm SFR}(z)$ has a pronounced peak at $z\sim1$--$2$ from the cosmic history of star formation, $w(z)+1$ is predicted to be {\it non-monotonic}: rising from $z=0$ to a peak at $z\sim1$--$2$, then falling back toward $w=-1$ at higher redshift (Fig.~\ref{fig:wz}).

This is fundamentally different from the CPL parameterization $w(a)=w_0+w_a(1-a)$, which is linear in $(1-a)$ and can only produce monotonic $w(z)$. No choice of CPL parameters can generate a peak. Quintessence models with simple potentials $V(\phi)\propto\phi^{-\alpha}$ also produce monotonic $w(z)$ \cite{tsujikawa2013}. The emergent dark energy model \cite{li2024} can produce non-monotonic behavior but lacks a connection to information theory.

\begin{figure}[htbp]
\centering
\includegraphics[width=\textwidth]{../figures/fig2_wz_shape.png}
\caption{\bf DGF $w(z)$ shape prediction.}
{\bf (a)} $w(z)$ for DGF with $n=1.0$ (blue), $n=0.7$ (green), and $n=1.3$ (red). The blue band shows the $1\sigma$ theory uncertainty from $\pm27\%$ SFR normalization and $\pm11\%$ cumulative stellar mass density. DESI $z\approx0$ constraint shown as red star. $\Lambda$CDM ($w\equiv-1$) shown as dotted line. The peak at $z\sim1$--$2$ is the key falsifiable prediction. {\bf (b)} $w(z)+1$ on logarithmic scale to highlight the peak structure. The shape at $z>2$ is sensitive to the memory kernel form (dashed continuation).
\label{fig:wz}
\end{figure}

Figure \ref{fig:wz} shows the DGF $w(z)$ prediction with its uncertainty band. The peak location and amplitude depend on $n$: larger $n$ shifts the peak to lower redshift and reduces its amplitude. For $n=1$, the peak is at $z\approx1.0$--$1.3$ with $w_{\rm peak}\approx-0.4$, a $\sim0.6$ deviation from $\Lambda$CDM. This is the ``smoking gun'' prediction: if DESI DR3 or Euclid data show $w(z)$ to be monotonic, DGF is falsified.

\section{Discussion}

We have shown that dark energy dynamics can be derived from two information-theoretic axioms, and that the resulting model agrees with DESI DR2 data within $1\sigma$---while being preferred over $\Lambda$CDM at $\sim7.5\sigma$. This is not curve-fitting: the shape function $w(z)+1\propto{\rm SFR}(z)/[H(z)\rho_*(z)^{n/(n+1)}]$ is a {\it consequence} of the physical mechanism (structure formation occupies vacuum information modes), not a parameterization chosen for its flexibility.

\subsection{Comparison with Competing Models}

Table \ref{tab:competitors} compares DGF against the main competing explanations for the DESI dynamical dark energy signal. DGF is unique in providing both a microphysical mechanism rooted in quantum information theory {\it and} independent falsifiable predictions beyond cosmology.

\begin{table}[htbp]
\centering
\caption{\bf Comparison of dark energy models.}
\label{tab:competitors}
\begin{tabular}{lccccc}
\toprule
Model & Physical mechanism? & Non-cosmology tests? & Non-monotonic $w(z)$? & $N_{\rm par}$ & $\chi^2$ (DESI DR2) \\
\midrule
$\Lambda$CDM & No & No & No & 0 & 57.5 \\
CPL & No (phenomenological) & No & No & 2 & 0 \\
Quintessence & Yes (scalar field) & No & No & 2--4 & $\sim$0 \\
Emergent DE \cite{li2024} & Partial & No & Possible & 2--3 & $\sim$0 \\
QMM \cite{neukart2025} & Yes (quantum memory) & In principle & No & 2 & $>100$ (excluded) \\
{\bf DGF (this work)} & {\bf Yes (information capacity)} & {\bf Yes (S1, S4, S5)} & {\bf Yes (peak)} & {\bf 2} & {\bf 0.65} \\
\bottomrule
\end{tabular}
\end{table}

\subsection{Falsifiable Predictions}

DGF makes three predictions that distinguish it from all competing models:

\begin{enumerate}
\item {\bf Non-monotonic $w(z)$ with a peak at $z\sim1$--$2$.} This is a ``smoking gun'' testable by DESI DR3 (2026--2027) and Euclid (2026+). If binned $w(z)$ measurements from these surveys are consistent with a monotonic function, DGF is falsified. No other model makes this prediction.

\item {\bf The specific shape $w(z)+1\propto{\rm SFR}(z)/[H(z)\rho_*(z)^{n/(n+1)}]$.} Once $n$ is constrained by low-$z$ data, this shape is a zero-additional-parameter prediction. Its agreement or disagreement with redshift-binned $w(z)$ measurements at $z>1$ will provide a decisive test.

\item {\bf Laboratory verification of the $q$-field.} In DGF, $q$ is not merely a proxy for $w(z)$---it has an independent operational meaning: the fraction of quantum information modes available for redundancy in quantum Darwinism \cite{zurek2009}. The DGF prediction $R_\delta(q)=R_\delta(1)\cdot q$ (where $R_\delta$ is the quantum Darwinism redundancy) can be tested in a $\sim$20-qubit superconducting circuit experiment within 2--3 years (see Supplementary \S6). This would give $q$ an operational definition completely independent of cosmology.
\end{enumerate}

\subsection{Honest Limitations}

We acknowledge several limitations of the current analysis:

\begin{enumerate}
\item {\bf CPL projection ambiguity.} The non-monotonic nature of DGF's $w(z)$ means that the CPL $(w_0,w_a)$ projection depends on the fitting range and weighting scheme. The true test of DGF is direct $w(z)$ shape comparison in redshift bins, not CPL parameters. Our CPL projection uses the DESI-sensitive range $z\in[0,2.3]$ with BAO volume weighting; different choices shift $w_a$ by $\pm0.15$.

\item {\bf SFR systematics dominate the theory error.} The $\pm27\%$ uncertainty in the local SFR normalization \cite{madau2014} contributes $\pm0.055$ to the $w_0+1$ error budget, exceeding DESI's measurement uncertainty of $\pm0.047$. Improved SFR measurements (e.g., H$\alpha$/UV/FUV multi-band) could reduce this to $\pm0.024$.

\item {\bf Memory kernel form.} We assume a constant memory kernel $K(\tau)\approx H_0c^2/R_c^2$. A frequency-dependent kernel (physically corresponding to different ``relaxation times'' for different structure formation modes) would modify the high-$z$ shape. This is simultaneously a limitation and a research opportunity: measuring $K(\tau)$ would probe the microphysics of information-mode relaxation.

\item {\bf $q$-field laboratory test not yet performed.} The quantum Darwinism experiment (Supplementary \S6) that would give $q$ an independent operational definition is currently a proposal. Until completed, DGF's cosmological interpretation of $q$ should be considered a model with an identified but unexecuted independent verification pathway.
\end{enumerate}

\section{Methods}

\subsection{DGF equation of state computation}

We solve the DGF telegraph equation (1) in the slow-roll regime coupled to the Friedmann equation $H^2(a)=H_0^2[\Omega_m a^{-3}+\Omega_r a^{-4}+\Omega_\Lambda q(a)]$. The coupling constant $\eta/\gamma_0$ is calibrated from the observed $w_0$. The memory kernel strength $\kappa$ is set by the memory-drive ratio $\mathcal{R}_0\equiv\kappa\int_0^{t_0}\Delta\,dt'/(\eta\,{\rm SFR}_0)$, which we take as $\mathcal{R}_0=0.3\pm0.2$ from requiring $w_0>-1$ (quintessence today).

\subsection{Star formation history}

We use the Madau \& Dickinson (2014) \cite{madau2014} cosmic star formation rate density:
\begin{equation}
{\rm SFR}(z)=0.015\,\frac{(1+z)^{2.7}}{1+[(1+z)/2.9]^{5.6}}\;M_\odot\,{\rm yr}^{-1}\,{\rm Mpc}^{-3}
\end{equation}
The cumulative stellar mass density is $\rho_*(z)=0.72\int_z^\infty{\rm SFR}(z')|dt/dz'|dz'$, where the factor 0.72 accounts for $\sim28\%$ mass return (Chabrier IMF) \cite{madau2014}. Uncertainties are $\pm27\%$ (SFR$_0$) and $\pm11\%$ ($\rho_{*,0}$) from Driver et al.\ (2018) \cite{driver2018}.

\subsection{Cosmological parameters}

We adopt Planck 2018 \cite{planck2018} values: $H_0=67.4$ km s$^{-1}$ Mpc$^{-1}$, $\Omega_m=0.315$, $\Omega_\Lambda=0.685$.

\subsection{DESI DR2 constraints}

We use the DESI+CMB+DESY5 joint constraints from Zhang et al.\ (2026) \cite{zhang2026}: $w_0=-0.785\pm0.047$, $w_a=-0.43\pm0.10$, with correlation coefficient $\rho=0.315$. The covariance matrix is:
\begin{equation}
\Sigma = \begin{pmatrix} 0.002209 & 0.001481 \\ 0.001481 & 0.0100 \end{pmatrix}
\end{equation}

\subsection{CPL projection}

We project DGF $w(z)$ onto the CPL parameterization $w(a)=w_0+w_a(1-a)$ using weighted least squares over the range $a\in[0.3,1]$ ($z\in[0,2.33]$), with DESI BAO sensitivity weights $\propto(1+z)^2/H(z)$. The $\chi^2$ is computed using the full covariance matrix.

\subsection{Code availability}

The complete numerical code, all figure data, and the machine-readable contour package are available at [repository URL]. The code is written in Python 3.12 using NumPy, SciPy, and Matplotlib.

\begin{thebibliography}{99}

\bibitem{riess1998}
Riess, A.G.\ et al.\ Observational evidence from supernovae for an accelerating universe and a cosmological constant. {\it Astron.\ J.} {\bf 116}, 1009--1038 (1998).

\bibitem{spergel2003}
Spergel, D.N.\ et al.\ First-year Wilkinson Microwave Anisotropy Probe (WMAP) observations: Determination of cosmological parameters. {\it Astrophys.\ J.\ Suppl.} {\bf 148}, 175--194 (2003).

\bibitem{tsujikawa2013}
Tsujikawa, S.\ Quintessence: a review. {\it Class.\ Quant.\ Grav.} {\bf 30}, 214003 (2013).

\bibitem{li2024}
Li, X.\ et al.\ Emergent dark energy: theoretical motivations and observational constraints. {\it Phys.\ Rev.\ D} {\bf 109}, 043510 (2024).

\bibitem{chevallier2001}
Chevallier, M.\ \& Polarski, D.\ Accelerating universes with scaling dark matter. {\it Int.\ J.\ Mod.\ Phys.\ D} {\bf 10}, 213--223 (2001).

\bibitem{linder2009}
Linder, E.V.\ \& Scherrer, R.J.\ Aetherizing Lambda: Barotropic fluids as dark energy. {\it Phys.\ Rev.\ D} {\bf 80}, 023008 (2009).

\bibitem{desi2025}
DESI Collaboration. DESI DR2 Results II: BAO measurements and cosmological constraints. arXiv:2503.14738 (2025).

\bibitem{zhang2026}
Zhang, X.\ et al.\ Robust evidence for dynamical dark energy from DESI DR2, CMB, and supernovae. SSRN 6215384 (2026).

\bibitem{madau2014}
Madau, P.\ \& Dickinson, M.\ Cosmic star-formation history. {\it Annu.\ Rev.\ Astron.\ Astrophys.} {\bf 52}, 415--486 (2014).

\bibitem{planck2018}
Planck Collaboration. Planck 2018 results. VI. Cosmological parameters. {\it Astron.\ Astrophys.} {\bf 641}, A6 (2020).

\bibitem{driver2018}
Driver, S.P.\ et al.\ GAMA/G10-COSMOS/3D-HST: the $0<z<5$ cosmic star formation history, stellar-mass, and dust-mass densities. {\it Mon.\ Not.\ R.\ Astron.\ Soc.} {\bf 475}, 2891--2935 (2018).

\bibitem{zurek2009}
Zurek, W.H.\ Quantum Darwinism. {\it Nature Phys.} {\bf 5}, 181--188 (2009).

\bibitem{neukart2025}
Neukart, F.\ et al.\ Quantum memory matrix model for dark energy. {\it Astronomy} {\bf 4}, 16 (2025).

\end{thebibliography}

\end{document}
